\section{Funtores representáveis}
\label{section:0.8}

\subsection{Funtores representáveis}
\label{subsection:0.8.1}

\oldpage[0\textsubscript{III}]{5}
\begin{env}[8.1.1]
\label{0.8.1.1}
Denotamos por $\Set$ a categoria de conjuntos.
Seja $\C$ uma categoria; para dois objetos $X$, $Y$ de $\C$, pomos $h_X(Y)=\Hom(Y,X)$; para cada morfismo $u:Y\to Y'$ de $\C$, denotamos por $h_X(u)$ a aplicação $v\mapsto vu$ de $\Hom(Y',X)$ em $\Hom(Y,X)$.
É imediato que, com estas definições, $h_X:\C\to\Set$ é um \emph{funtor contravariante}, isto é, um objeto da categoria $\CHom(\C\op,\Set)$ dos funtores covariantes da categoria $\C\op$ (a categoria dual de $\C$) na categoria $\Set$ (T, 1.7, (d) e \cite{III-29}).
\end{env}

\begin{env}[8.1.2]
\label{0.8.1.2}
Seja agora $w:X\to X'$ um morfismo de $\C$; para cada $Y\in\C$ e cada $v\in\Hom(Y,X)=h_X(Y)$,
temos $wv\in\Hom(Y,X')=h_{X'}(Y)$; denotamos por $h_w(Y)$ a aplicação $v\mapsto wv$ de
$h_X(Y)$ em $h_{X'}(Y)$. É imediato que, para cada morfismo $u:Y\to Y'$ de $\C$, o
diagrama
\[
  \xymatrix{
    h_X(Y')\ar[r]^{h_X(u)}\ar[d]_{h_w(Y')} &
    h_X(Y)\ar[d]^{h_w(Y)}\\
    h_{X'}(Y')\ar[r]_{h_{X'}(u)} &
    h_{X'}(Y)
  }
\]
é comutativo; em outros termos, $h_w$ é uma \emph{transformação natural (ou morfismo funtorial)
$h_X\to h_{X'}$} (T, 1.2), e também um morfismo da categoria $\CHom(\C\op,\Set)$ (T, 1.7, (d)). As
definições de $h_X$ e de $h_w$ constituem, portanto, a definição de um \emph{funtor covariante
canônico}
\[
  h:\C\to\CHom(\C\op,\Set),\quad X\mapsto h_X.
  \tag{8.1.2.1}
\]
\end{env}

\begin{env}[8.1.3]
\label{0.8.1.3}
Seja $X$ um objeto de $\C$, e seja $F$ um funtor contravariante de $\C$ em $\Set$
(um objeto de $\CHom(\C\op,\Set)$). Seja $g:h_X\to F$ uma \emph{transformação natural}: para todo $Y\in\C$,
\oldpage[0\textsubscript{III}]{6}
$g(Y)$ é, pois, uma aplicação $h_X(Y)\to F(Y)$ tal que, para cada morfismo $u:Y\to Y'$ de $\C$,
o diagrama
\[
  \xymatrix{
    h_X(Y')\ar[r]^{h_X(u)}\ar[d]_{g(Y')} &
    h_X(Y)\ar[d]^{g(Y)}\\
    F(Y')\ar[r]_{F(u)} &
    F(Y)
  }
  \tag{8.1.3.1}\label{0.8.1.3.1}
\]
é comutativo. Em particular, temos uma aplicação $g(X):h_X(X)=\Hom(X,X)\to F(X)$ e, por conseguinte, um elemento
\[
  \alpha(g)=(g(X))(1_X)\in F(X)
  \tag{8.1.3.2}
\]
e, consequentemente, uma aplicação canônica
\[
  \alpha:\Hom(h_X,F)\to F(X).
  \tag{8.1.3.3}
\]

Reciprocamente, consideremos um elemento $\xi\in F(X)$; para cada morfismo $v:Y\to X$ de $\C$, $F(v)$ é uma
aplicação $F(X)\to F(Y)$; consideremos a aplicação
\[
  v\mapsto(F(v))(\xi)
  \tag{8.1.3.4}
\]
de $h_X(Y)$ em $F(Y)$; se denotarmos esta aplicação por $(\beta(\xi))(Y)$,
\[
  \beta(\xi):h_X\to F
  \tag{8.1.3.5}
\]
é uma \emph{transformação natural}, pois para cada morfismo $u:Y\to Y'$ de $\C$ temos
$(F(vu))(\xi)=(F(u)\circ F(v))(\xi)$, o que torna (\hyperref[0.8.1.3.1]{8.1.3.1}) comutativo para $g=\beta(\xi)$.
Temos definido assim uma aplicação canônica
\[
  \beta:F(X)\to\Hom(h_X,F).
  \tag{8.1.3.6}
\]
\end{env}

\begin{proposition}[8.1.4]
\label{0.8.1.4}
As aplicações $\alpha$ e $\beta$ são bijeções inversas uma da outra.
\end{proposition}

\begin{proof}
Calculemos $\alpha(\beta(\xi))$ para $\xi\in F(X)$; para cada $Y\in\C$, $(\beta(\xi))(Y)$ é a aplicação
$g_1(Y):v\mapsto(F(v))(\xi)$ de $h_X(Y)$ em $F(Y)$. Temos, pois,
\[
  \alpha(\beta(\xi))=(g_1(X))(1_X)=(F(1_X))(\xi)=1_{F(X)}(\xi)=\xi.
\]
Calculemos agora $\beta(\alpha(g))$ para $g\in\Hom(h_X,F)$; para cada $Y\in\C$, $(\beta(\alpha(g)))(Y)$
é a aplicação $v\mapsto(F(v))((g(X))(1_X))$; pela comutatividade de (\hyperref[0.8.1.3.1]{8.1.3.1}), esta aplicação
não é senão $v\mapsto(g(Y))((h_X(v))(1_X))=(g(Y))(v)$, por definição de $h_X(v)$; em outros termos,
é igual a $g(Y)$, o que termina a demonstração.
\end{proof}

\begin{env}[8.1.5]
\label{0.8.1.5}
Recordemos que uma \emph{subcategoria} $\C'$ de uma categoria $\C$ se define pela condição de que seus objetos
sejam objetos de $\C$, e de que, se $X'$, $Y'$ são dois objetos de $\C'$, o conjunto
$\Hom_{\C'}(X',Y')$ de morfismos $X'\to Y'$ \emph{em $\C'$} seja um subconjunto do conjunto $\Hom_\C(X',Y')$ de
morfismos $X'\to Y'$ \emph{em $\C$}, sendo a aplicação canônica de ``composição de morfismos''
\[
  \Hom_{\C'}(X',Y')\times\Hom_{\C'}(Y',Z')\to\Hom_{\C'}(X',Z')
\]
\oldpage[0\textsubscript{III}]{7}
a restrição da aplicação canônica
\[
  \Hom_\C(X',Y')\times\Hom_\C(Y',Z')\to\Hom_\C(X',Z').
\]

Dizemos que $\C'$ é uma subcategoria \emph{plena} de $\C$ se $\Hom_{\C'}(X',Y')=\Hom_\C(X',Y')$ para todo
par de objetos de $\C'$. A subcategoria $\C''$ de $\C$ formada pelos objetos de $\C$ isomorfos a
objetos de $\C'$ é então também uma subcategoria plena de $\C$, \emph{equivalente} (T, 1.2) a $\C'$, como se
verifica facilmente.

Um funtor covariante $F:\C_1\to\C_2$ chama-se \emph{plenamente fiel} se, para todo par de objetos
$X_1$, $Y_1$ de $\C_1$, a aplicação $u\mapsto F(u)$ de $\Hom(X_1,Y_1)$ em $\Hom(F(X_1),F(Y_1))$ é
\emph{bijetiva}; isto implica que a subcategoria $F(\C_1)$ de $\C_2$ é \emph{plena}. Além disso,
se dois objetos $X_1$, $X_1'$ têm a mesma imagem $X_2$, existe um único isomorfismo
$u:X_1\to X_1'$ tal que $F(u)=1_{X_2}$. Para cada objeto $X_2$ de $F(\C_1)$, seja $G(X_2)$ um dos
os objetos $X_1$ de $\C_1$ tais que $F(X_1)=X_2$ ($G$ se define mediante o axioma da escolha); para
cada morfismo $v:X_2\to Y_2$ de $F(\C_1)$, $G(v)$ será o único morfismo $u:G(X_2)\to G(Y_2)$ tal
que $F(u)=v$; $G$ é então um \emph{funtor} de $F(\C_1)$ em $\C_1$; $FG$ é o funtor identidade de
$F(\C_1)$, e o anterior mostra que existe um isomorfismo de funtores $\vphi:1_{\C_1}\to GF$
tal que $F$, $G$, $\vphi$ e a identidade $1_{F(\C_1)}\to FG$ definem uma \emph{equivalência} entre
a categoria $\C_1$ e a subcategoria plena $F(\C_1)$ de $\C_2$ (T, 1.2).
\end{env}

\begin{env}[8.1.6]
\label{0.8.1.6}
Aplicamos a Proposição \sref{0.8.1.4} ao caso em que $F$ é $h_{X'}$, sendo $X'$ um
objeto qualquer de $\C$; a aplicação $\beta:\Hom(X,X')\to\Hom(h_X,h_{X'})$ não é senão a aplicação
$w\mapsto h_w$ definida em \sref{0.8.1.2}; uma vez que esta aplicação é \emph{bijetiva}, vemos,
com a terminologia de \sref{0.8.1.5}, que:
\end{env}

\begin{proposition}[8.1.7]
\label{0.8.1.7}
O funtor canônico $h:\C\to\CHom(\C\op,\Set)$ é plenamente fiel.
\end{proposition}

\begin{env}[8.1.8]
\label{0.8.1.8}
Seja $F$ um funtor contravariante de $\C$ em $\Set$; dizemos que $F$ é \emph{representável} se
existe um objeto $X\in\C$ tal que $F$ é \emph{isomorfo} a $h_X$; da
Proposição \sref{0.8.1.7} segue-se que os dados de um $X\in\C$ e de um isomorfismo de funtores
$g:h_X\to F$ determinam $X$ salvo isomorfismo único. A Proposição \sref{0.8.1.7} implica então
que $h$ define uma \emph{equivalência} entre $\C$ e a subcategoria plena de
$\CHom(\C\op,\Set)$ formada pelos \emph{funtores contravariantes representáveis}. Da
Proposição \sref{0.8.1.4} segue-se que o dado de uma transformação natural $g:h_X\to F$ é
equivalente ao de um elemento $\xi\in F(X)$; dizer que $g$ é um \emph{isomorfismo} equivale
à seguinte condição sobre $\xi$: \emph{para todo objeto $Y$ de $\C$, a aplicação
$v\mapsto(F(v))(\xi)$ de $\Hom(Y,X)$ em $F(Y)$ é bijetiva}. Quando $\xi$ satisfaz esta condição,
dizemos que o par $(X,\xi)$ \emph{representa} o funtor representável $F$. Por abuso de linguagem,
dizemos também que o objeto $X\in\C$ representa $F$ se existe um $\xi\in F(X)$ tal que
$(X,\xi)$ representa $F$, isto é, se $h_X$ é isomorfo a $F$.

Sejam $F$, $F'$ dois funtores contravariantes representáveis de $\C$ em $\Set$, e sejam $h_X\to F$ e
$h_{X'}\to F'$ dois isomorfismos de funtores. Segue-se então de \sref{0.8.1.6} que
existe uma correspondência bijetiva canônica entre $\Hom(X,X')$ e o conjunto $\Hom(F,F')$ de
transformações naturais $F\to F'$.
\end{env}

\begin{env}[8.1.9]
\label{0.8.1.9}
\textit{Exemplo I. Limites projetivos}.
A noção de funtor contravariante representável compreende, em particular, a noção ``dual'' da noção usual de ``solução de um problema universal''.
\oldpage[0\textsubscript{III}]{8}
Mais geralmente, veremos que a noção de \emph{limite projetivo} é um caso particular da noção de funtor representável.
Recordemos que, em uma categoria $\C$, define-se um \emph{sistema projetivo} pelos dados de um conjunto pré-ordenado $I$, de uma família $(A_\alpha)_{\alpha\in I}$ de objetos de $\C$ e, para cada par de índices $(\alpha,\beta)$ tal que $\alpha\leq\beta$, de um morfismo $u_{\alpha\beta}:A_\beta\to A_\alpha$.
Um \emph{limite projetivo} deste sistema em $\C$ consiste em um objeto $B$ de $\C$ (denotado $\varprojlim A_\alpha$) e, para cada $\alpha\in I$, um morfismo $u_\alpha:B\to A_\alpha$ tais que: 1\textsuperscript{o}. $u_\alpha=u_{\alpha\beta}u_\beta$ para $\alpha\leq\beta$; 2\textsuperscript{o}. para todo objeto $X$ de $\C$ e toda família $(v_\alpha)_{\alpha\in I}$ de morfismos $v_\alpha:X\to A_\alpha$ tal que $v_\alpha=u_{\alpha\beta}v_\beta$ para $\alpha\leq\beta$, existe um único morfismo $v:X\to B$ (denotado $\varprojlim v_\alpha$) tal que $v_\alpha=u_\alpha v$ para todo $\alpha\in I$ (T, 1.8).
Isto pode ser interpretado da seguinte maneira: os $u_{\alpha\beta}$ definem canonicamente aplicações
\[
  \overline{u}_{\alpha\beta}:\Hom(X,A_\beta)\to\Hom(X,A_\alpha)
\]
que definem um \emph{sistema projetivo} de conjuntos $(\Hom(X,A_\alpha),\overline{u}_{\alpha\beta})$, e $(v_\alpha)$ é, por definição, um elemento do conjunto $\varprojlim\Hom(X,A_\alpha)$; é claro que $X\mapsto\varprojlim\Hom(X,A_\alpha)$ é um \emph{funtor contravariante} de $\C$ em $\Set$, e a existência do limite projetivo $B$ equivale a dizer que $(v_\alpha)\mapsto\varprojlim v_\alpha$ é um \emph{isomorfismo} de funtores em $X$
\[
  \varprojlim\Hom(X,A_\alpha)\isoto\Hom(X,B),
  \tag{8.1.9.1}
\]
isto é, que o funtor $X\mapsto\varprojlim\Hom(X,A_\alpha)$ é \emph{representável}.
\end{env}

\begin{env}[8.1.10]
\label{0.8.1.10}
\textit{Exemplo II. Objetos finais}.
Seja $\C$ uma categoria, e seja $\{a\}$ um conjunto unitário.
Consideremos o funtor contravariante $F:\C\to\Set$ que envia cada objeto $X$ de $\C$ ao conjunto $\{a\}$, e cada morfismo $X\to X'$ de $\C$ à única aplicação $\{a\}\to\{a\}$.
Dizer que este funtor é \emph{representável} significa que existe um objeto $e\in\C$ tal que, para todo $Y\in\C$, $\Hom(Y,e)=h_e(Y)$ é um \emph{conjunto unitário}; dizemos que $e$ é um \emph{objeto final} de $\C$, e é claro que dois objetos finais de $\C$ são isomorfos (o que permite definir, em geral mediante o axioma da escolha, \emph{um} objeto final de $\C$, que denotamos então por $e_\C$).
Por exemplo, na categoria $\Set$, os objetos finais são os conjuntos unitários; na categoria de \emph{álgebras aumentadas} sobre um corpo $K$ (onde os morfismos são os homomorfismos de álgebras compatíveis com o aumento), $K$ é um objeto final; na categoria de \emph{$S$-esquemas} \sref[I]{I.2.5.1}, $S$ é um objeto final.
\end{env}

\begin{env}[8.1.11]
\label{0.8.1.11}
Para dois objetos $X$ e $Y$ de uma categoria $\C$, ponhamos $h_X'(Y)=\Hom(X,Y)$ e, para todo morfismo $u:Y\to Y'$, seja $h_X'(u)$ a aplicação $v\mapsto uv$ de $\Hom(X,Y)$ em $\Hom(X,Y')$; $h_X'$ é então um \emph{funtor covariante $\C\to\Set$}, de modo que deduzimos, como em \sref{0.8.1.2}, a definição de um funtor covariante canônico $h':\C\op\to\CHom(\C,\Set)$; um funtor \emph{covariante} $F$ de $\C$ em $\Set$, isto é, um objeto de $\CHom(\C,\Set)$, é então \emph{representável} se existe um objeto $X\in\C$ (necessariamente único salvo isomorfismo único) tal que $F$ é \emph{isomorfo a $h_X'$}; deixamos ao leitor o desenvolvimento das noções ``duais'' das anteriores, que desta vez compreendem a noção de \emph{limite indutivo}, e em particular a noção usual de ``solução de um problema universal''.
\end{env}

\subsection{Estruturas algébricas em categorias}
\label{subsection:0.8.2}

\begin{env}[8.2.1]
\label{0.8.2.1}
\oldpage[0\textsubscript{III}]{9}
Dados dois funtores contravariantes $F$ e $F'$ de $\C$ em $\Set$, recordemos que, para todo objeto $Y\in\C$, pomos $(F\times F')(Y)=F(Y)\times F'(Y)$, e para todo morfismo $u:Y\to Y'$ de $\C$, pomos $(F\times F')(u)=F(u)\times F'(u)$, que é a aplicação $(t,t')\mapsto(F(u)(t),F'(u)(t'))$ de $F(Y')\times F'(Y')$ em $F(Y)\times F'(Y)$; $F\times F':\C\to\Set$ é assim um \emph{funtor contravariante} (que não é senão o \emph{produto} dos objetos $F$ e $F'$ na categoria $\CHom(\C\op,\Set)$).
Dado um objeto $X\in\C$, chamamos \emph{lei de composição interna} sobre $X$ a uma \emph{transformação natural}
\[
\label{0.8.2.1.1}
  \gamma_X:h_X\times h_X\to h_X.
  \tag{8.2.1.1}
\]

Em outros termos (T, 1.2), para todo objeto $Y\in\C$, $\gamma_X(Y)$ é uma aplicação $h_X(Y)\times h_X(Y)\to h_X(Y)$ (portanto, por definição, uma \emph{lei de composição interna} sobre o conjunto $h_X(Y)$) com a condição de que, para todo morfismo $u:Y\to Y'$ de $\C$, o diagrama
\[
  \xymatrix{
    h_X(Y')\times h_X(Y')\ar[rr]^{h_X(u)\times h_X(u)}\ar[dd]_{\gamma_X(Y')} & &
    h_X(Y)\times h_X(Y)\ar[dd]^{\gamma_X(Y)}\\\\
    h_X(Y')\ar[rr]_{h_X(u)} & &
    h_X(Y)
  }
\]
seja comutativo; isto implica que, para as leis de composição $\gamma_X(Y)$ e $\gamma_X(Y')$, $h_X(u)$ é um \emph{homomorfismo} de $h_X(Y')$ em $h_X(Y)$.

De maneira análoga, dados dois objetos $Z$ e $X$ de $\C$, chamamos \emph{lei de composição externa sobre $X$, com $Z$ como domínio de operadores} a uma transformação natural
\[
\label{0.8.2.1.2}
  \omega_{X,Z}:h_Z\times h_X\to h_X.
  \tag{8.2.1.2}
\]

Vemos, como antes, que, para todo $Y\in\C$, $\omega_{X,Z}(Y)$ é uma lei de composição externa sobre $h_X(Y)$, com $h_Z(Y)$ como domínio de operadores, e tal que, para todo morfismo $u:Y\to Y'$, $h_X(u)$ e $h_Z(u)$ formam um \emph{di-homomorfismo} de $(h_Z(Y'),h_X(Y'))$ em $(h_Z(Y),h_X(Y))$.
\end{env}

\begin{env}[8.2.2]
\label{0.8.2.2}
Seja $X'$ um segundo objeto de $\C$, e suponhamos dada uma lei de composição interna $\gamma_{X'}$ sobre $X'$; dizemos que um morfismo $w:X\to X'$ de $\C$ é um \emph{homomorfismo} para as leis de composição se, para todo $Y\in\C$, $h_w(Y):h_X(Y)\to h_{X'}(Y)$ é um \emph{homomorfismo} para as leis de composição $\gamma_X(Y)$ e $\gamma_{X'}(Y)$.
Se $X''$ é um terceiro
\oldpage[0\textsubscript{III}]{10}
objeto de $\C$ munido de uma lei de composição interna $\gamma_{X''}$, e $w':X'\to X''$ é um morfismo de $\C$ que é um homomorfismo para $\gamma_{X'}$ e $\gamma_{X''}$, é claro que o morfismo $w'w:X\to X''$ é um homomorfismo para as leis de composição $\gamma_X$ e $\gamma_{X''}$.
Um isomorfismo $w:X\isoto X'$ de $\C$ chama-se \emph{isomorfismo para as leis de composição $\gamma_X$ e $\gamma_{X'}$} se $w$ é um homomorfismo para estas leis de composição, e se seu morfismo inverso $w^{-1}$ é um homomorfismo para as leis de composição $\gamma_{X'}$ e $\gamma_X$.

Definimos de maneira análoga os \emph{di-homomorfismos} para pares de objetos de $\C$ munidos de leis de composição externa.
\end{env}

\begin{env}[8.2.3]
\label{0.8.2.3}
Quando uma lei de composição interna $\gamma_X$ sobre um objeto $X\in\C$ é tal que $\gamma_X(Y)$ é uma \emph{lei de grupo} sobre $h_X(Y)$ para \emph{todo} $Y\in\C$, dizemos que $X$, munido desta lei, é um \emph{$\C$-grupo} ou um \emph{objeto grupo de $\C$}.
Definimos analogamente os \emph{$\C$-anéis}, \emph{$\C$-módulos}, etc.
\end{env}

\begin{env}[8.2.4]
\label{0.8.2.4}
Suponhamos que o \emph{produto} $X\times X$ de um objeto $X\in\C$ por si mesmo existe em $\C$; por definição, temos então $h_{X\times X}=h_X\times h_X$ salvo isomorfismo canônico, pois é um caso particular do limite projetivo \sref{0.8.1.9}; uma lei de composição interna sobre $X$ pode, portanto, ser considerada como um morfismo funtorial $\gamma_X:h_{X\times X}\to h_X$, e determina assim canonicamente \sref{0.8.1.6} um elemento $c_X\in\Hom(X\times X,X)$ tal que $h_{c_X}=\gamma_X$; neste caso, o dado de uma lei de composição interna sobre $X$ equivale ao dado de um morfismo $X\times X\to X$; quando $\C$ é a categoria $\Set$, recuperamos a noção clássica de lei de composição interna sobre um conjunto.
Temos um resultado análogo para uma lei de composição externa quando o produto $Z\times X$ existe em $\C$.
\end{env}

\begin{env}[8.2.5]
\label{0.8.2.5}
Com a notação anterior, suponhamos além disso que $X\times X\times X$ existe em $\C$; a caracterização do produto como objeto que representa um funtor \sref{0.8.1.9} implica a existência de isomorfismos canônicos
\[
  (X\times X)\times X\isoto X\times X\times X\isoto X\times(X\times X);
\]
se identificarmos canonicamente $X\times X\times X$ com $(X\times X)\times X$, a aplicação $\gamma_X(Y)\times 1_{h_X(Y)}$ identifica-se com $h_{c_X\times 1_X}(Y)$ para todo $Y\in\C$.
Consequentemente, é equivalente dizer que, para todo $Y\in\C$, a lei interna $\gamma_X(Y)$ é associativa, que o diagrama de aplicações
\[
  \xymatrix{
    h_X(Y)\times h_X(Y)\times h_X(Y)\ar[rr]^{\gamma_X(Y)\times 1}\ar[dd]_{1\times\gamma_X(Y)} & &
    h_X(Y)\times h_X(Y)\ar[dd]^{\gamma_X(Y)}\\\\
    h_X(Y)\times h_X(Y)\ar[rr]_{\gamma_X(Y)} & &
    h_X(Y)
  }
\]
\oldpage[0\textsubscript{III}]{11}
é comutativo, e que o diagrama de morfismos
\[
  \xymatrix{
    X\times X\times X\ar[rr]^{c_X\times 1_X}\ar[dd]_{1_X\times c_X} & &
    X\times X\ar[dd]^{c_X}\\\\
    X\times X\ar[rr]_{c_X} & &
    X
  }
\]
é comutativo.
\end{env}

\begin{env}[8.2.6]
\label{0.8.2.6}
Sob as hipóteses de \sref{0.8.2.5}, se quisermos expressar, para todo $Y\in\C$, que a lei interna $\gamma_X(Y)$ é uma \emph{lei de grupo}, é necessário, em primeiro lugar, que seja associativa e, em segundo lugar, que exista uma aplicação $\alpha_X(Y):h_X(Y)\to h_X(Y)$ com as propriedades da operação de \emph{inverso} de um grupo; como vimos que, para todo morfismo $u:Y\to Y'$ de $\C$, $h_X(u)$ deve ser um homomorfismo de grupos $h_X(Y')\to h_X(Y)$, vemos primeiro que $\alpha_X:h_X\to h_X$ deve ser uma \emph{transformação natural}.
Por outro lado, podem-se expressar as propriedades características do inverso $s\mapsto s^{-1}$ de um grupo $G$ sem fazer intervir o elemento identidade: basta verificar que as duas aplicações compostas
\[
  (s,t)\mapsto(s,s^{-1},t)\mapsto(s,s^{-1}t)\mapsto s(s^{-1}t),
\]
\[
  (s,t)\mapsto(s,s^{-1},t)\mapsto(s,ts^{-1})\mapsto(ts^{-1})s
\]
são iguais à segunda projeção $(s,t)\mapsto t$ de $G\times G$ em $G$.
Por \sref{0.8.1.3}, temos $\alpha_X=h_{a_X}$, onde $a_X\in\Hom(X,X)$; a primeira condição anterior expressa então que o morfismo composto
\[
  X\times X\xrightarrow{(1_X,a_X)\times 1_X}X\times X\times X\xrightarrow{1_X\times c_X}X\times X\xrightarrow{c_X}X
\]
é a segunda projeção $X\times X\to X$ em $\C$, e a segunda condição é análoga.
\end{env}

\begin{env}[8.2.7]
\label{0.8.2.7}
Suponhamos agora que existe um \emph{objeto final} $e$ \sref{0.8.1.10} em $\C$.
Suponhamos sempre que $\gamma_X(Y)$ é uma lei de grupo sobre $h_X(Y)$ para todo $Y\in\C$, e denotemos por $\eta_X(Y)$ o elemento identidade de $\gamma_X(Y)$.
Como, para todo morfismo $u:Y\to Y'$ de $\C$, $h_X(u)$ é um homomorfismo de grupos, temos $\eta_X(Y)=(h_X(u))(\eta_X(Y'))$; tomando em particular $Y'=e$, caso no qual $u$ é o único elemento $\varepsilon$ de $\Hom(Y,e)$, vemos que o elemento $\eta_X(e)$ determina por completo $\eta_X(Y)$ para todo $Y\in\C$.
Ponhamos $e_X=\eta_X(X)$, o elemento identidade do grupo $h_X(X)=\Hom(X,X)$; a comutatividade do diagrama
\[
  \xymatrix{
    h_X(e)\ar[r]^{h_X(\varepsilon)}\ar[d]_{h_{e_X}(e)} &
    h_X(Y)\ar[d]^{h_{e_X}(Y)}\\
    h_X(e)\ar[r]_{h_X(\varepsilon)} &
    h_X(Y)
  }
\]
(cf. \sref{0.8.1.2}) mostra que, sobre o conjunto $h_X(Y)$, a aplicação $h_{e_X}(Y)$ não é senão
\oldpage[0\textsubscript{III}]{12}
$s\mapsto\eta_X(Y)$, que envia todo elemento ao elemento identidade.
Comprovamos então que o fato de $\eta_X(Y)$ ser o elemento identidade de $\gamma_X(Y)$ para todo $Y\in\C$ equivale a dizer que o morfismo composto
\[
  X\xrightarrow{(1_X,1_X)}X\times X\xrightarrow{1_X\times e_X}X\times X\xrightarrow{c_X}X,
\]
e seu análogo no qual se intercambiam $1_X$ e $e_X$, são ambos \emph{iguais} a $1_X$.
\end{env}

\begin{env}[8.2.8]
\label{0.8.2.8}
Naturalmente, poder-se-iam ampliar com facilidade os exemplos de estruturas algébricas em categorias.
O exemplo dos grupos foi tratado com suficiente detalhe, mas daqui em diante deixaremos, em geral, ao leitor o desenvolvimento das noções análogas para os exemplos de estruturas algébricas que encontrarmos.
\end{env}
